\documentclass[UTF8,a4paper,10.5pt]{ctexart}
\usepackage[left=1.9cm,right=1.9cm,top=1.8cm,bottom=1.8cm]{geometry}
\usepackage{amsmath,amssymb}
\usepackage{enumitem}
\usepackage{fancyhdr}
\usepackage{lastpage}
\usepackage{listings}
\usepackage{xcolor}
\usepackage{array}

\pagestyle{fancy}
\fancyhf{}
\renewcommand{\headrulewidth}{0pt}
\fancyfoot[C]{第~\thepage~页\quad 共~\pageref{LastPage}~页}

\lstset{
    language=C,
    basicstyle=\ttfamily\small,
    keywordstyle=\color{blue!70!black}\bfseries,
    commentstyle=\color{green!45!black},
    stringstyle=\color{red!60!black},
    columns=fullflexible,
    keepspaces=true,
    breaklines=true,
    showstringspaces=false,
    escapeinside={(*@}{@*)},
    aboveskip=2pt,
    belowskip=2pt
}

\lstdefinestyle{outputstyle}{
    language={},
    basicstyle=\ttfamily\small,
    keywordstyle={},
    commentstyle={},
    stringstyle={},
    columns=fullflexible,
    keepspaces=true,
    breaklines=true,
    showstringspaces=false,
    aboveskip=2pt,
    belowskip=4pt,
    frame=leftline,
    framerule=1.5pt,
    rulecolor=\color{gray!50},
    xleftmargin=1.2em
}

\newcommand{\blankans}[1]{\underline{\ \texttt{#1}\ }}
\newcommand{\parenval}[1]{\hfill(\makebox[1.4cm][c]{#1})}

\begin{document}

\begin{center}
    {\Large\bfseries 2020 级信息学院《C 语言程序设计》考试试题（A）参考答案}
\end{center}

\medskip

\noindent\textbf{一、判断下列语句或程序段的对错。（“$\times$”表示错，“$\checkmark$”表示对）（10分）}

\begin{enumerate}[label=(\arabic*), leftmargin=2.2em, itemsep=2pt]
    \item \texttt{int x=05678;} \parenval{$\times$}
    \item \texttt{int x = 0x177f;}\\
          \texttt{float y = 1.8e+5;} \parenval{$\checkmark$}
    \item \lstinline|char _a01='\101';| \parenval{$\checkmark$}
    \item \texttt{int **p, (*pa)[10];}\\
          \texttt{p = pa;} \parenval{$\times$}
    \item \texttt{int (*pa)[5],a[5][5];}\\
          \texttt{pa = a++;} \parenval{$\times$}
    \item \texttt{char *str[10];}\\
          \texttt{gets(str[0]);} \parenval{$\times$}
    \item \texttt{int *p, a[5][5];}\\
          \texttt{p = *a+1;} \parenval{$\checkmark$}
    \item \texttt{int a[5]=\{0, 1, 2,3,4,5\};} \parenval{$\times$}
    \item \texttt{char **p, *pa[10];}\\
          \texttt{p = pa;} \parenval{$\checkmark$}
    \item \lstinline|char *str = "\\tabcd\123\\n";| \parenval{$\checkmark$}
\end{enumerate}

\medskip

\noindent\textbf{二、计算下列表达式的值（10分）}

\noindent 假设 \texttt{int} 和 \texttt{unsigned int} 类型均为 16 位长度，且各题彼此独立。

\noindent 设 \texttt{unsigned int a=2, b=3; int c=4, d=5; float f; char str[10]="hust";}
\begin{enumerate}[label=(\arabic*), leftmargin=2.2em, itemsep=2pt]
    \item \texttt{f = c/b + d/a;} \parenval{\texttt{3.0}}
    \item \texttt{--a || --a \&\& a;} \parenval{\texttt{1}}
    \item \texttt{(a\^{}c) \& (b|d)} \parenval{\texttt{6}}
    \item \texttt{f = sizeof(str) + strlen(str);} \parenval{\texttt{14.0}}
    \item \texttt{a-c > d - b ? a : b;} \parenval{\texttt{2}}
\end{enumerate}

\newpage
\noindent\textbf{三、程序改错（10分）要求：不得改变程序框架，不得重写程序，无需文字说明，直接在代码上添加、删除和修改。}

\noindent 1、如下程序对已知的 \texttt{N} 个字符串排序，其中在 \texttt{main()} 函数中初始化，在 \texttt{sort} 函数中排序（使用指针数组）。
\begin{lstlisting}
#include <stdio.h>
#include <string.h>               // 添加头文件
#define N 5                       // 添加符号常量定义
void sort(char *pstr[]);          // 添加函数声明

void main()
{
    int i;
    char str[N][80] = {"basic", "pascal", "colbol", "losp", "agol"};
    char *pstr[N] = {str[0], str[1], str[2], str[3], str[4]};  // 修改：初始化指针数组
    sort(pstr);
    for (i = 0; i < N; i++)
    {
        printf("%s\n", pstr[i]);  // 修改：输出 pstr[i]
    }
}

void sort(char *pstr[])           // 修改：形参改为 char *pstr[]
{
    int i, j;
    char *temp;
    for (i = 0; i < N - 1; i++)       // 修改：循环条件改为 i < N - 1
        for (j = 0; j < N - 1 - i; j++)  // 修改：循环条件改为 j < N - 1 - i
            if (strcmp(pstr[j], pstr[j + 1]) > 0)  // 修改：用 strcmp 比较字符串大小
            {
                temp = pstr[j];
                pstr[j] = pstr[j + 1];
                pstr[j + 1] = temp;
            }
}
\end{lstlisting}

\medskip
\noindent 2、设有 4 个候选人，\texttt{N} 个人参加选举，每次输入一个得票的候选人的名字，要求最后输出个人的得票结果。
\begin{lstlisting}
#include <stdio.h>
#include <string.h>               // 添加头文件
#define N 10
struct person
{
    char name[10];                // 修改：成员改为字符数组（或 char *name;）
    int count;
};

void main()
{
    struct person leader[4] = {{"wang", 0},
                               {"zhao", 0},
                               {"mao",  0},
                               {"gao",  0}};
    int i, j;
    char name[10];                // 添加字符数组定义
    for (i = 0; i < N; i++)
    {
        gets(name);
        for (j = 0; j < 4; j++)
            if (strcmp(leader[j].name, name) == 0)  // 修改：用 strcmp 比较且下标改为 j
            {
                leader[j].count++;  // 修改：对 leader[j].count 自增
                break;              // 修改：改为 break
            }
    }
    printf("\n");
    for (j = 0; j < 4; j++)
        printf("%s:%d\n", leader[j].name, leader[j].count);
}
\end{lstlisting}

\medskip

\noindent\textbf{四、程序填空（10分）}

\noindent (1)~求已知两个正整数的最大公约数。
\begin{lstlisting}
#include <stdio.h>

void main()
{
    (*@\blankans{int i}@*);
    int a, b;

    printf("please input a ,b:");
    scanf("%d%d", (*@\blankans{\&a, \&b}@*));

    for (i = (*@\blankans{a>b?b:a}@*); i > 0; i--)  // i 初值为 a,b 中的较小值
    {
        if ((*@\blankans{a\%i==0 \&\& b\%i==0}@*))
        {
            printf("the max  = %d ", i);
            (*@\blankans{break}@*);
        }
    }
}
\end{lstlisting}

\medskip
\noindent (2)~求序列：$1!,\ 2!,\ 3!,\ 4!\cdots\cdots$ 的前 10 项之和。
\begin{lstlisting}
(*@\blankans{\#include <stdio.h>}@*)

void main()
{
    int i;
    long int sum, t;  // sum 代表和, t 代表某项

    sum = 0;
    t = 1;

    for (i = 1; (*@\blankans{i<=10}@*); i++)
    {
        sum = (*@\blankans{sum+t}@*);
        (*@\blankans{t = t*(i+1)}@*);
    }
    printf(" sum = %ld", (*@\blankans{sum}@*));
}
\end{lstlisting}

\medskip
\noindent\textbf{五、输出程序运行结果（25分，每题5分，结果写在题目的右边）}

\noindent 1、
\begin{lstlisting}
#include <stdio.h>

void main()
{
    int i;
    for (i = 1; i < 10; i++)
    {
        if (i % 3)
        {
            putchar('W');
        }
        else if (i % 6 == 0) break;
        else continue;
        printf("U  HAN\n");
    }
    printf("I  love you");
}
\end{lstlisting}
\noindent\textbf{运行结果：}
\begin{lstlisting}[style=outputstyle]
WU  HAN
WU  HAN
WU  HAN
WU  HAN
I  love you
\end{lstlisting}

\medskip
\noindent 2、
\begin{lstlisting}
#include <stdio.h>

int n;

void display()
{
    int a = 1;
    static int b = 1;
    printf("a=%d,  b=%d  c=%d\n", a++, b++, ++n);
}

void main()
{
    int i;
    for (i = 0; i < 5; i++)
    {
        display();
    }
}
\end{lstlisting}
\noindent\textbf{运行结果：}
\begin{lstlisting}[style=outputstyle]
a=1,  b=1  c=1
a=1,  b=2  c=2
a=1,  b=3  c=3
a=1,  b=4  c=4
a=1,  b=5  c=5
\end{lstlisting}

\medskip
\noindent 3、
\begin{lstlisting}
#include <stdio.h>

void main()
{
    char str[][20] = {"china",
                      "hubei",
                      "wuhan",
                      "hust-aia",
                      "c language"};
    int i;
    for (i = 0; i < 5; i++)
    {
        printf("%s\n", str[4 - i] + 4 - i);
    }
}
\end{lstlisting}
\noindent\textbf{运行结果：}
\begin{lstlisting}[style=outputstyle]
nguage
t-aia
han
ubei
china
\end{lstlisting}

\medskip
\noindent 4、
\begin{lstlisting}
#include <stdio.h>

void main()
{
    char *str[] = {"pascal", "cobol", "fortran", "lisp", "java"};
    char **p[] = {str + 4, str + 3, str + 2, str + 1, str};
    char ***pp = p;

    printf("%s---", **(pp++) + 1);
    printf("%s---", **(pp++) + 2);
    printf("%s---", **(pp++) + 1);
    printf("%s",    **(pp++) + 2);
    printf("---%s", **(pp));
}
\end{lstlisting}
\noindent\textbf{运行结果：}
\begin{lstlisting}[style=outputstyle]
ava---sp---ortran---bol---pascal
\end{lstlisting}

\medskip
\noindent 5、
\begin{lstlisting}
#include <stdio.h>
#define N 6

void main()
{
    int a[N][N];
    int i, j;

    for (i = 0; i < N; i++)
    {
        a[i][0] = 1;
        a[i][i] = 1;
    }
    for (i = 2; i < N; i++)
    {
        for (j = 1; j < i; j++)
        {
            a[i][j] = a[i - 1][j - 1] + a[i - 1][j];
        }
    }
    for (i = 0; i < N; i++)
    {
        for (j = 0; j <= i; j++)
        {
            printf("%5d", a[i][j]);
        }
        printf("\n");
    }
}
\end{lstlisting}
\noindent\textbf{运行结果：}
\begin{lstlisting}[style=outputstyle]
    1
    1    1
    1    2    1
    1    3    3    1
    1    4    6    4    1
    1    5   10   10    5    1
\end{lstlisting}

\medskip

\noindent\textbf{六、编程（35分）注意：不得使用全局变量，注意程序结构}

\noindent 1、(1)~编一程序，输入 $x$ 的值，输出 $y$ 的值。（6分）
\[
y = \begin{cases}
x + 10, & x < 10 \\
2x + 200, & 10 \le x < 100 \\
3x + 300, & x \ge 100
\end{cases}
\]

\noindent 2、编写一猜数游戏程序，随机产生某个小于 $N$ 的正整数，从键盘反复输入整数进行猜数，当未猜中时，提示输入过大或过小。猜中时，指出猜的次数。最多允许猜 20 次。（8分）

\noindent 3、编一程序，要求主函数中输入一行英文和一个字母，被调用的函数返回删除该字母后的那行英文（如有多个，一并删除）并在主函数中输出。（9分）

\noindent 4、某班有学生 $N$ 名，每名学生信息由姓名、性别、学号、成绩 1、成绩 2 和成绩 3 组成，试编程要求：
\begin{enumerate}[label=（\arabic*）, leftmargin=2.4em, itemsep=2pt]
    \item 学生信息由键盘输入
    \item 求每人平均分
    \item 按平均分数从高到低排序
    \item 按平均分从高到低显示所有平均分不及格的男生信息。
\end{enumerate}
\noindent （每小题写一函数，通过 \texttt{main} 函数调用实现）（12分）

\medskip
\noindent\textbf{参考程序：}

\medskip
\noindent 1、
\begin{lstlisting}
#include <stdio.h>

void main()
{
    float x, y;
    scanf("%f", &x);
    if (x < 10)
    {
        y = x + 10;
    }
    else if (x < 100)
    {
        y = 2 * x + 200;
    }
    else
    {
        y = 3 * x + 300;
    }
    printf("y=%f", y);
}
\end{lstlisting}

\medskip
\noindent 2、
\begin{lstlisting}
#include <stdio.h>
#include <time.h>
#include <stdlib.h>
#define N 100

void main()
{
    int a, i, x, count;
    srand((unsigned)time(NULL));
    a = rand() % (N - 1) + 1;
    count = 0;
    for (i = 0; i < 20; i++)
    {
        scanf("%d", &x);
        count++;
        if (x == a)
        {
            break;
        }
        else if (x < a)
        {
            printf("太小");
        }
        else
        {
            printf("太大");
        }
    }
    if (i < 20)
    {
        printf("猜中用的次数=%d", count);
    }
    else
    {
        printf("未猜中");
    }
}
\end{lstlisting}

\medskip
\noindent 3、
\begin{lstlisting}
#include <stdio.h>
#include <string.h>

char *delchar(char *pstr, char c);

void main()
{
    char str[80], *ps;
    char c;
    gets(str);
    scanf("%c", &c);
    ps = delchar(str, c);
    puts(ps);
}

char *delchar(char *pstr, char c)
{
    static char p[80];
    int i, j;
    for (i = 0, j = 0; pstr[i] != '\0'; i++)
    {
        if (pstr[i] != c)
        {
            p[j++] = pstr[i];
        }
    }
    p[j] = '\0';
    return p;
}
\end{lstlisting}

\medskip
\noindent 4、
\begin{lstlisting}
#include <stdio.h>
#define N 30

struct student
{
    char name[10];
    char sex;
    long int no;
    float score[3];
    float ave;
};

void input(struct student *p);
void average(struct student *p);
void sort(struct student *p);
void findnopassm(struct student *p);

void main()
{
    struct student stu[N];
    input(stu);
    average(stu);
    sort(stu);
    findnopassm(stu);
}

void input(struct student *p)
{
    int i;
    for (i = 0; i < N; i++)
    {
        scanf("%s %c%ld%f%f%f", p[i].name, &p[i].sex, &p[i].no,
              &p[i].score[0], &p[i].score[1], &p[i].score[2]);
    }
}

void average(struct student *p)
{
    int i;
    for (i = 0; i < N; i++)
    {
        p[i].ave = (p[i].score[0] + p[i].score[1] + p[i].score[2]) / 3;
        printf("%s\t%ld\t%f\n", p[i].name, p[i].no, p[i].ave);
    }
}

void sort(struct student *p)
{
    struct student temp;
    int i, j;
    for (i = 0; i < N - 1; i++)
        for (j = 0; j < N - 1 - i; j++)
        {
            if (p[j].ave < p[j + 1].ave)
            {
                temp = p[j];
                p[j] = p[j + 1];
                p[j + 1] = temp;
            }
        }
    for (i = 0; i < N; i++)
    {
        printf("%s\t%ld\t%f\n", p[i].name, p[i].no, p[i].ave);
    }
}

void findnopassm(struct student *p)
{
    int i;
    for (i = 0; i < N; i++)
        if (p[i].sex == 'M' && p[i].ave < 60)
        {
            printf("%s\t%c\t%ld\t%f\n", p[i].name, p[i].sex, p[i].no, p[i].ave);
        }
}
\end{lstlisting}

\end{document}
